\documentclass[11pt]{article}
\usepackage[utf8]{inputenc}
\usepackage{amsmath}
\usepackage{mathtools}
\usepackage{lipsum} % to generate some filler text
\usepackage{fullpage}
\usepackage{mathrsfs}
\usepackage{svg}
\usepackage{xr}
\usepackage{amsthm,cite,url,graphicx,booktabs,lipsum,color,bm,caption,subcaption,soul}
\usepackage{pifont,tikz,paralist,multirow,amssymb}
\usepackage{xifthen}
\usepackage{enumerate}
\usepackage{titlesec}
\newcounter{reviewer}
\setcounter{reviewer}{0}
\newcounter{point}[reviewer]
\setcounter{point}{0}

\usepackage{hyperref}

% This refines the format of how the reviewer/point reference will appear.
%\renewcommand{\thepoint}{C\,\thereviewer.\arabic{point}}

\renewcommand{\thepoint}{Comment\,\thereviewer.\arabic{point}}

% command declarations for reviewer points and our responses

\newenvironment{point1}
   {\refstepcounter{point} \bigskip \noindent {\textbf{Associate Editor Comment~}} ---\ }
   {\par }

\newcommand{\reviewersection}{\stepcounter{reviewer} \bigskip \hrule
                  \section*{Reviewer \thereviewer}}

% \newenvironment{point}

%    {\refstepcounter{point} \bigskip \noindent {\textbf{Reviewer~\thepoint} } ---\ }
%    {\par }

\newenvironment{point}
   {\refstepcounter{point} \bigskip \noindent {\textbf{Reviewer~\thepoint} } ---\ }
   {\par }


\newcommand{\shortpoint}[1]{\refstepcounter{point}   \bigskip \noindent 
	{\textbf{Reviewer~Point~\thepoint} } ---~#1  }

\newenvironment{reply}
   {\medskip \noindent \color{blue} \begin{sf}\textbf{Reply}:\  }
   {\medskip \par \end{sf}}

\newcommand{\shortreply}[2][]{\medskip \noindent \begin{sf}\textbf{Reply}:\  #2
	\ifthenelse{\equal{#1}{}}{}{ \hfill \footnotesize (#1)}%
	\medskip \end{sf}}

\makeatletter
\newcommand*{\addFileDependency}[1]{% argument=file name and extension
  \typeout{(#1)}
  \@addtofilelist{#1}
  \IfFileExists{#1}{}{\typeout{No file #1.}}
}
\makeatother

\newcommand*{\myexternaldocument}[1]{%
    \externaldocument{#1}%
    \addFileDependency{#1.tex}%
    \addFileDependency{#1.aux}%
}

\def\thesubsubsection{\arabic{subsubsection}}
\titleformat{\subsubsection}[runin]
{\normalfont\bfseries}
{\indent\thesubsubsection)}{0.5em}{}

% FOR LOCAL RUNS ONLY: COMMENT FOR OVERLEAF
% Resolve \ref to labels in the revised manuscript (requires sn-article-revised.aux)
\myexternaldocument{sn-article-revised}


\begin{document}
\section*{Responses to reviewer comments on \textit{Explainable Binary Classification of Separable Shape Ensembles}}
\vspace{10pt}

The authors would like to thank the editor and reviewers for their constructive feedback and suggestions that have greatly helped to improve the quality of this manuscript. %We have carefully revised the manuscript by addressing all the comments of reviewers. Additionally, we have paraphrased certain sections of the manuscript without altering the original meaning and interpretation to address copyright considerations relating to our conference paper~\cite{Set} that was published at the same time as this review. This conference paper has also been introduced in section:~\textquotedblleft Article contributions" and we additionally have made a note of the unique contributions in this article that expands upon that previous work. The ~\textquotedblleft Conclusions" section has been revised to include additional details based on the comments from the reviewers. For reviewers' convenience, we have used latexdiff to mark up all insertions, deletions, and modifications in the revised manuscript.

Please see below our responses to reviewer's comments highlighted in \textcolor{blue}{blue}.

% Let's start point-by-point with Reviewer 1

\section*{Response to the reviewers}

%%%%%%%%%%%%%%%%%%%%%%%%%%%%%%%%%%%%%%%%%%%%%%%%%%%
\reviewersection

The authors propose a method for statistical comparison of shape samples based on segmented images. The main idea is to approximate a segmented curve via a (mathematical) landmark configuration and then decompose variation into an object that captures undulations (nonlinearity) and another object that captures generalized scales. The former is an element of the Grassmann manifold while the latter is an element of the space of symmetric, positive-definite matrices. A cyclic Procrustes algorithm is introduced to account for arbitrary rotations of and starting points on the landmark configurations. After the decomposition, the authors perform principal component analysis (PCA) in appropriate tangent spaces to reduce dimension and represent the objects via Euclidean features (PCs). Finally, statistical comparison of the shape ensembles is enabled via a maximum mean discrepancy (MMD)-based hypothesis test. The framework is applied to compare electron backscatter diffraction images of materials samples.

I read this paper with interest and feel that it presents a valuable addition to the statistical shape analysis literature. However, there are some issues as I outline next that should be addressed in a revision.
%%%%%%%%%%%%%%%%

\begin{point}
Sections 2 and 3. While generally the presentation of concepts was clear and easy to understand, there are portions of these two sections that were cumbersome to read due to heavy notation and lack of illustrations. I encourage the authors to revise these two sections by (i) simplifying notation in certain places, and (ii) providing illustrations of some of the main results.
\end{point}
\begin{reply}
We have taken several concrete steps to address readability of Sections~\ref{sec:FIE}--\ref{sec:numerics}. (i)~Lemmas~\ref{lemma:symm}, \ref{lemma:HS_int_opt}, and \ref{lemma:CLO_invariance}, together with Theorem~\ref{theorem:eigfuncs}, have been rewritten in the standard ``Let\,$\ldots$\,Then$\ldots$'' style. (ii)~The composite linear operator and PRRTI condition are now presented in numbered definition environments (Definitions~\ref{def:CLO} and~\ref{def:PRRTI}). (iii)~Two new illustrations have been introduced: Figure~\ref{fig:SRQD} (Section~\ref{subsec:func_approx}) visualizes the SRQD as a weighted polar decomposition, and Figure~\ref{fig:placeholder} (Section~\ref{subsec:mfld_learn}) gives a pipeline diagram mapping preshapes through registration to normal coordinates over the product submanifold. (iv)~Much of the lengthier SST background has been relocated to Appendix~\ref{secA0} (``Review of Separable Shape Tensors''), and the retained in-text review (Section~\ref{subsec:SST}) has been renamed ``Review'' accordingly. A further schematic contrasting nonlinear undulation with linear scale deformations is in preparation.
\end{reply}

\begin{point}
Terminology. I feel that the title and terminology used in the paper is a bit misleading. The authors are not performing classification, but rather a statistical comparison via a hypothesis testing procedure based on MMD. Further, in Table 1 and Section 3.5, I would change the terminology used for the decision made after hypothesis testing from “Accept” to “Fail to Reject Null”.
\end{point}
\begin{reply}
We respectfully disagree that the terminology is misleading: two-sample hypothesis testing defines a binary decision rule mapping pairs of distributions to $\{$reject, FTR$\}$, which is itself a binary classification---albeit distinct from supervised classification over labeled instances. We have retained the title and added supporting material. The abstract has been revised to foreground the hypothesis-testing framing. In Section~\ref{sec:intro}, three new blue-text paragraphs formally introduce the binary mapping $H:(\mathcal{E},\widehat{\mathcal{E}})\mapsto\{0,1\}$, its separable decomposition $H(\mathcal{E},\widehat{\mathcal{E}}) \coloneqq H(\mathcal{G}_r,\widehat{\mathcal{G}}_r)\wedge H(\mathcal{P},\widehat{\mathcal{P}})$, and its extension to context-driven classification tasks. Section~\ref{subsec:MMD} mirrors this notation in blue text connecting pMMD to the binary mapping and explicitly distinguishing this unsupervised, distribution-level classification from supervised learning over labeled instances. Per your suggestion, Table~\ref{tab:logical_conjunction} and all adjacent text in Sections~\ref{subsec:MMD} and~\ref{sec:experiments} now use ``FTR'' in place of ``Accept.''
\end{reply}

\begin{point}
Cyclic Procrustes. The authors use a cyclic Procrustes algorithm to pre-register the landmark configurations with respect to rotation and the seed (starting point of the configuration). For this, all shapes within an image are aligned to a randomly selected template configuration. In my opinion, a randomly selected template is not a very good choice to perform this registration. The result can be quite sensitive to this choice, especially if an outlying shape is selected as a template. Why didn’t the authors perform this registration with respect to the Frechet mean of the ensemble? While I understand that this would increase computational cost, it would result in more stable registration results.
\end{point}
\begin{reply}
The reviewer's suggestion is well-motivated. A blue-text justification has been added in Section~\ref{subsec:cycl_Pro} following Theorem~\ref{thm:cycl_Procrustes}, making three points. (i)~Theorem~\ref{thm:cycl_Procrustes} guarantees \emph{an} optimal rotation and cyclic permutation aligning each curve to \emph{any} asymmetric archetype; the essential requirement is asymmetry of the template, not its optimality as a Fr\'echet mean. (ii)~Fr\'echet-mean alignment is a chicken-and-egg problem---the mean presupposes aligned data while alignment presupposes a template---and an iterative scheme lacks convergence guarantees for the discrete cyclic-permutation component; iteratively re-solving $N$ SVDs across $n$ quadrature nodes becomes expensive for large ensembles. (iii)~Empirical stability of the random-archetype registration is demonstrated by the decision-landscape experiments in Figures~\ref{fig:010_decision} and~\ref{fig:010_decision-normalized}: consistent decisions across different quadrature resolutions $n$ (which correspond to different discrete realizations of the same random archetype) and even under random row-wise permutations provide strong evidence that the registration is not pathologically sensitive to the draw.
\end{reply}

\begin{point}
Details of MMD-based Hypothesis Test. How was the p-value computed? It is my understanding that the MMD-based hypothesis test is a permutation test, because the distribution of the test statistic is not given in closed form and must be estimated empirically. There is a lack of details of how the permutation test was implemented: (i) how many permutations were used? (ii) how was the cyclic Procrustes procedure incorporated into the hypothesis test? In particular, question (ii) raises an important issue. Since a permutation of the shapes changes the ensembles, cyclic Procrustes should be carried out each time. I did not see a description of this in the manuscript.
\end{point}
\begin{reply}
We agree that the procedural detail was insufficient, and a blue-text expansion to Section~\ref{subsec:MMD} is in preparation specifying: (i)~the unbiased $\widehat{\mathrm{MMD}}^2$ estimator from Lemma~6 of Gretton et al.\ (2012), computed separately for undulation and scale coordinates; (ii)~a Gaussian RBF kernel with median pairwise-distance bandwidth; (iii)~$B=1000$ permutations over exchangeable ensemble labels (image~1 vs.\ image~2); and (iv)~the permutation p-value $p = (1 + \#\{b : \widehat{\mathrm{MMD}}^2_b \geq \widehat{\mathrm{MMD}}^2_{\mathrm{obs}}\})/(1+B)$, rejecting at level $\alpha$ if $p\leq\alpha$. Regarding your important point~(ii): cyclic Procrustes is performed \emph{once} against an asymmetric archetype using the \emph{pooled aggregate ensemble} \emph{before} the permutation loop. Because the archetype is drawn from the union of both images, the registration is label-agnostic; permutations therefore shuffle pre-registered normal coordinates $(\boldsymbol{t}_q,\boldsymbol{\ell}_q)$ rather than raw curves, which preserves validity of the permutation null while avoiding re-running $N{+}\widehat{N}$ registrations per permutation.
\end{reply}

\begin{point}
Simulations. To validate the overall framework, could the authors carry-out a study based on simulated curves/landmark configurations? This would further help assess the hypothesis testing procedure, e.g., type I error and power.
\end{point}
\begin{reply}
A simulated-curve study assessing Type~I error and power is in preparation. We also note that the decision-landscape sensitivity studies already included in Section~\ref{sec:experiments} (Figures~\ref{fig:010_decision} and~\ref{fig:010_decision-normalized}) offer partial empirical evidence: in the $(0\wedge 1)$ case, decisions are consistent at FTR for $r\leq 7$ and consistent at reject for $r\geq 8$ across many independent quadrature discretizations, bounding empirical error rates near the transition and providing an indirect stability analysis complementary to a controlled simulation.
\end{reply}

\begin{point}
The authors should provide an idea of the computational cost associated with the entire procedure.
\end{point}
\begin{reply}
A per-stage runtime table (interpolation/arc-length reparametrization, SRQD, cyclic Procrustes, tangent PCA, permutation MMD) on a standard laptop is in preparation and will appear in Section~\ref{sec:experiments}. We note that the two dominant costs are cyclic Procrustes (linear in the number of curves, searching a subset of $p$-cycles per shape) and the $B=1000$ permutation loop (each permutation recomputes $\widehat{\mathrm{MMD}}^2$ on pre-registered coordinates, so no registration is repeated). The full pipeline for the EBSD experiments in Section~\ref{sec:experiments} executes in seconds on a conventional laptop for thousands of curves.
\end{reply}

%%%%%%%%%%%%%%%%%%%%%%%%%%%%%%%%%%%%%%%%%%%%%%%%%%%
\reviewersection

This work presents a formal mathematical framework for representing discretized curve shape data by linear scale variations and nonlinear undulations using separable shape tensors, which is used to motivate the statistical comparison of shape ensembles between pairs of images based on maximum mean discrepancy (MMD).  I think this is a promising contribution, as rigorous comparisons of large shape collections is an important research area, and this appears to be a good fit for this journal with the mathematical rigor. I also particularly liked the numerical experiments in Section 4, which studied practical comparisons specifically in the context of image processing choices. However, I would suggest some modifications, described below.
%%%%%%%%%%%%%%%%

\begin{point}
Is it reasonable to assume separability of undulation and scale in real data examples?  Can they be correlated, and if so, how could this be diagnosed?  This seems to have strong ramifications with regards to later discussion of obtaining separate coordinates for each using tangent PCA, and the MMD-based hypothesis tests.
\end{point}
\begin{reply}
This is an important distinction, now clarified in blue text in Section~\ref{subsec:MMD}. \emph{Separability} refers to the algebraically exact polar decomposition $T_n = \widetilde{X}P$ valid for any full-rank $T_n$; it does not assume or imply statistical independence of the resulting normal coordinates $(\boldsymbol{t},\boldsymbol{\ell})$. When physical processes simultaneously drive grain growth and boundary roughness (e.g., recrystallization), joint variability may carry supplemental explanatory information. A diagnostic has also been added as blue text: one may test whether the joint distribution factors as the product of its marginals using $\mathrm{MMD}[\mathcal{F},\boldsymbol{\rho},\rho_{\boldsymbol{t}}\rho_{\boldsymbol{\ell}}]$ (equivalent to the Hilbert-Schmidt Independence Criterion). When independence fails, the pMMD and the De~Morgan decomposition in Table~\ref{tab:logical_conjunction} remain valid logical statements about marginal discrepancies, but a direct test of the joint distribution can augment the separable explanations.
\end{reply}

\begin{point}
I’m not sure if “explainable binary classification” is the right terminology here.  The authors present MMD to compare shape ensembles from pairs of images to assess if there are discrepancies (via hypothesis tests).  One could certainly use the normal coordinates applied from separate tangent PCA steps in Section 3.4 to build a classification model which takes an image and outputs a binary class label, but the task of interest is not in classification itself, but rather assessing if shape distributions between images are the same or not, and attributing any differences to scale or undulation variation. I would recommend a different title (e.g., “Explainable Discrepancy Tests of Separable Shape Ensembles”).
\end{point}
\begin{reply}
We have retained the title. As we now argue in blue text in Section~\ref{sec:intro} and Section~\ref{subsec:MMD}, two-sample hypothesis testing \emph{is} a binary classification: the map $H:(\mathcal{E},\widehat{\mathcal{E}})\mapsto\{$reject, FTR$\}$ is a binary decision rule over pairs of distributions. The qualifier ``explainable'' refers to the SST decomposition, which identifies \emph{why} a discrepancy arises---undulation, scale, or both---via the logical conjunction in Table~\ref{tab:logical_conjunction}. A binary-decision machine over pairs of ensembles also naturally generalizes to context-driven classification: holding a trusted ensemble fixed and varying an unknown ``query'' ensemble yields label assignment, as elaborated in the new blue-text paragraph at the end of Section~\ref{sec:intro}. We have also updated ``Accept'' to ``FTR'' throughout per reviewer R1.
\end{reply}

\begin{point}
While there is plenty of mathematical rigor, there is little statistical detail, particularly in Section 3. For example, it appears that the authors use a non-parametric two-sample hypothesis test with MMD as the fundamental quantity; but these details are not provided. Given that this is the focus of Section 3 and the experiments in Section 4, these details should be provided to help improve reader understanding.
\end{point}
\begin{reply}
We agree, and refer the reviewer to our response to R1.4. A blue-text expansion to Section~\ref{subsec:MMD} in preparation specifies the non-parametric permutation test in full: the unbiased $\widehat{\mathrm{MMD}}^2$ estimator from Lemma~6 of Gretton et al.\ (2012), Gaussian RBF kernel with median-distance bandwidth, $B=1000$ label permutations performed on pre-registered coordinates, and the conservative permutation p-value. Coupled with the separability clarification in response to R2.1 and the multiple-testing correction discussed in response to R2.4, Section~\ref{subsec:MMD} will contain a complete statistical specification of the procedure used in Section~\ref{sec:experiments}.
\end{reply}

\begin{point}
In Section 3.5, the authors discuss a singular test involving a product MMD metric, stating its equivalence to performing separate MMD-based tests involving undulation and generalized scales and using DeMorgan’s logic rules (Table 1). This doesn’t take into account multiple testing concerns: if the separate undulation MMD and generalized scale MMD tests both are constructed individually at level alpha, this does not mean the joint test has level alpha, and this can lead to more type I errors for the joint hypothesis, unless adjustments are made (e.g., Bonferroni corrections, etc.) to control the type I error for the joint test.  These details should be discussed further.
\end{point}
\begin{reply}
The reviewer raises an important point. A Bonferroni correction---conducting each marginal MMD test at level $\alpha/2$ so the joint Type~I error under the intersection null is bounded by $\alpha$---is being added as blue text to Section~\ref{subsec:MMD}. For our experiments at $\alpha=0.01$, each marginal test uses $\alpha/2=0.005$; this correction does not change any of the reported decisions in Section~\ref{sec:experiments} because all rejections satisfy $p\ll 0.005$ and all FTR cases satisfy $p\gg 0.01$. We also clarify the relationship to De~Morgan: the De~Morgan equivalence is a \emph{logical} identity relating the joint and marginal hypotheses (dictating \emph{which} tests to run), while Bonferroni is a \emph{statistical} adjustment to control the error rate of the logical conjunction (dictating at \emph{what level} to run them); the two are entirely compatible.
\end{reply}

\begin{point}
More discussion on selecting r, n, particularly within Section 4 (where these were fixed) would help.  Should these be chosen relatively large, or to achieve some threshold of cumulative variance explained?  This could have an impact on tests, since it is possible that discrepancies in shape ensembles between image pairs only show up at later dimensions (which contribute little to the overall cumulative variance).
\end{point}
\begin{reply}
Practical guidance for selecting $r$ (undulation dimensionality) and $n$ (quadrature nodes) is being added as blue text in Section~\ref{sec:experiments}. Our primary recommendation is to leverage the decision-landscape framework itself (Figures~\ref{fig:010_decision}, \ref{fig:010_decision-normalized}, and~\ref{fig:avg_decision_land}) as a diagnostic: (i)~$n$ should be large enough to visually resolve boundary features, with decisions observed to be stable for $n\in[100,500]$ thanks to the spectral convergence of the quadrature (Figure~\ref{fig:nystrom_eg}); (ii)~$r$ is the more consequential parameter---the reviewer is correct that discrepancies may appear only at higher dimensions, as illustrated in the $(0\wedge 1)$ case where decisions flip at $r=8$. We therefore recommend selecting $r$ via a cumulative-variance threshold (e.g.\ 95--99\%) from the tangent-PCA spectrum, then \emph{verifying stability} by scanning the decision landscape over a range of $r$ values and examining the ambiguity metric introduced in Section~\ref{subsec:MMD}; a choice beyond the ambiguity peak is preferred. A dedicated supplemental experiment formalizing these criteria is under consideration.
\end{reply}

\begin{point}
The writing was a bit hard to follow: I found the organization a little bit confusing, in that I could not distinguish which of these results were from existing work vs. new to this manuscript, and the language felt a bit cumbersome. Some suggestions:
\begin{itemize}
    \item A figure which describes nonlinear shape undulations vs. linear scale differences would be useful.
    \item The relative amounts of detail for different section doesn’t seem to reflect the contributions and focus of the work.  For example, the paper’s title suggests a focus on MMD and hypothesis testing/explaining discrepancies, but none of this is really discussed until page 17.  I would suggest moving some of the details in Section 2-3 to a supplementary section, if possible.
\end{itemize}
\end{point}
\begin{reply}
We have taken several steps toward clarifying contributions and tightening the organization. The bulk of the SST background has been relocated to Appendix~\ref{secA0} (``Review of Separable Shape Tensors''), with the in-text review condensed and renamed in Section~\ref{subsec:SST}. The contributions discussion in Section~\ref{sec:intro} has been updated to cite \cite{grey2023sst} at the point of primary novelty and to delineate new material from prior work. New blue-text comparisons contrast our approach with varifold-based methods~\cite{kaltenmark2017,pierson2022}, Sobolev metrics on curves~\cite{bauer2017}, and AI/deep-learning approaches~\cite{hartman2023,hartman2021}. We plan a further concision pass and an undulation-versus-scale schematic in the next revision.

Regarding wholesale relocation of Sections~2--3 to supplementary material: we respectfully prefer to keep the core formalisms in the main narrative. We believe hiding mathematical details in supplementary material undermines transparency and reproducibility, especially for a mathematical journal---a concern we share with the classical tradition of weaving formalism into the narrative. We will continue to pursue concision within the main body.
\end{reply}

\begin{point}
Typo on page 6 after Eq. (5): ``it sufficient'' $\rightarrow$ ``it is sufficient''
\end{point}
\begin{reply}
Fixed (Section~\ref{subsec:preshapes}, following Eq.~\eqref{eq:reg_Hilbert_space}). A full copy-editing pass also addressed ``elude''$\to$``allude'' (two occurrences), ``eigendecompsition''$\to$``eigendecomposition'', ``apriori''$\to$``\emph{a priori}'', ``nickle''$\to$``nickel'', doubled ``of of'' and ``by by'', several subject--verb agreement errors, and two reference errors. See CHANGES.md for the complete log.
\end{reply}


%%%%%%%%%%%%%%%%%%%%%%%%%%%%%%%%%%%%%%%%%%%%%%%%%%%
\reviewersection
This paper presents a novel method for ensembles of curves leveraging so called separable shape tensors (SST) (introduced in \cite{grey2023sst}) arising from eigenspaces of composite integral operators. This approach offers rigid-invariant decompositions of curves into generalized scale variations and undulations. The use of SST representations of curves allows for binary classification of samples of curves utilizing a product maximum mean discrepancy with interpretable feature spaces allowing for explainable binary classification. This work presents a robust set of experiments highlighting the effectiveness and explainability of their model with focus on comparing distributions of curves
segmented from images. 

This paper has several significant strengths:
\begin{itemize}
    \item The numerical experiments support the paper nicely and support the paper’s approach to binary shape classification and offer clear evidence of of the explainability of the model.
    \item The paper offers a lot of math background to the approach to explainable binary classification of curves. The proofs are presented with sufficient detail for the reader to follow. My primary feedback is that the presentation of the statements of these results can be hard to parse, but the proofs themselves are well presented.
\end{itemize}
Recommendation: Accept after revisions, focusing on improving the formal presentation of the statements of Theorems, Lemmas, and Definitions of key concepts.

%%%%%%%%%%%%%%%%
\begin{point}
The statements of lemmas and theorems in Section 2 are presented with varying levels of
formality. Specifically:
\begin{itemize}
    \item[a.] Lemma 1, Lemma 3, and Theorem 1 could be presented more formally. i.e. for Lemma 1 ``Let $c \in \mathcal{C}_d$, then ...''
    \item[b.] Lemma 2 reads more formally, but the wording of ``$\mathcal{K} \tau [c]$ such that $\mathcal{K} \tau [c] = \dots$ is awkward and the operator could just be introduced as $\mathcal{K} \tau [c] \coloneqq \dots$.
\end{itemize}
\end{point}
\begin{reply}
Done. Lemma~\ref{lemma:symm}, Lemma~\ref{lemma:CLO_invariance}, and Theorem~\ref{theorem:eigfuncs} have been rewritten in the standard ``Let\,$\ldots$\,Then$\ldots$'' form with explicit hypotheses. Lemma~\ref{lemma:HS_int_opt} has been reworded to introduce the Hilbert-Schmidt integral operator via $\coloneqq$ notation as suggested, with ``Then $\mathcal{K}_{\mathcal{T}}[\boldsymbol{c}]$ is a Hilbert-Schmidt (HS) integral operator.'' stated as the conclusion.
\end{reply}

\begin{point}
Along with this, I found Section 2 and Section 3 felt difficult to follow due to the use of acronyms. Especially when they are used in the context of mathematical statements. For instance, PRRTI and CLO have rigorous mathematical definitions and would benefit from being presented in a definition environment. This would make it easier for a reader unfamiliar with these terms to quickly reference the meaning of these terms rather than needing to look for the first in-text uses of these terms. In particular,
\begin{itemize}
    \item[a.] I suggest using a definition environment for the CLO or dropping the acronym entirely and using the symbol for the composite linear operator throughout. It is a bit confusing having an acronym and math symbol being used interchangeably throughout the paper.
    \item[b.] For the PRRTI, I would suggest ``Definition N (PRRTI condition). We say an operator is permutation, rotation, rescale, and translation invariant (PRRTI) if for any $c \in \mathcal{C}_d$, $R \in O(d)$, ...''
\end{itemize}
\end{point}
\begin{reply}
Done. The composite linear operator is now introduced in Definition~\ref{def:CLO} (Section~\ref{sec:FIE}), and the PRRTI condition in Definition~\ref{def:PRRTI}, using essentially the phrasing you suggested (``permutation, rotation, reflection, and translation invariant''). We have retained the acronyms because they are referenced throughout subsequent formal statements (e.g., Theorem~\ref{theorem:eigfuncs} and Lemma~\ref{lemma:CLO_invariance}), but each usage is now anchored to its definition environment. Definition~\ref{def:eval_func} is retained for the RKHS evaluation functionals from~\cite{grey2023sst}.
\end{reply}

\begin{point}
The paper begins using acronyms such as SST for separable shape tensor and HS for Hilbert-Schmidt without introducing the acronym on the initial use. This makes the paper difficult to follow at times.
\end{point}
\begin{reply}
Fixed. SST (separable shape tensor) is now introduced parenthetically at its first occurrence in the abstract and reinforced at the first in-text use in Section~\ref{sec:intro} and Section~\ref{subsec:SST}. HS (Hilbert-Schmidt) is introduced parenthetically at its first use inside Lemma~\ref{lemma:HS_int_opt} (``\ldots\ is a Hilbert-Schmidt (HS) integral operator.''). Other acronyms---CLO, PRRTI, tPCA, SRQD, pMMD, FTR---are introduced similarly on first use.
\end{reply}

\begin{point}
The paper contains several typos, difficult to parse sentences, and incorrect word choices i.e.
”elude” vs. ”allude”. The paper would benefit from a pass for grammar.
\end{point}
\begin{reply}
A full copy-editing pass has been completed. Corrections include: ``elude''$\to$``allude'' at both occurrences (Sections~\ref{sec:FIE} and~\ref{sec:e_solns}), ``it sufficient''$\to$``it is sufficient'' (Section~\ref{subsec:preshapes}), ``eigendecompsition''$\to$``eigendecomposition'' (Section~\ref{subsec:geo_interp}), ``apriori''$\to$``\emph{a priori}'' (Section~\ref{subsec:cycl_Pro}), ``nickle''$\to$``nickel'' (Section~\ref{sec:experiments}), ``weighted version of of''$\to$``weighted version of'', ``by by taking''$\to$``by taking'', ``These numerical experiment emphasizes''$\to$``These numerical experiments emphasize'', ``SST's''$\to$``SSTs'', a missing \texttt{Figure\textasciitilde} prefix, and a misused \texttt{\textbackslash eqref} in Section~\ref{subsec:Nystrom}. Several difficult-to-parse sentences in the introduction were simplified. See CHANGES.md for the complete log.
\end{reply}

%%%%%%%%%%%%%%%%%%%%%%%%%%%%%%%%%%%%%%%%%%%%%%%%%%%
\reviewersection
The paper proposes a pipeline for classification of curves, applied to curves ensemble extracted from EBSD (electron microscopy imaging). The pipeline is, to my understanding as follow: 
\begin{itemize}
    \item Embed curves as landmarks in Rnxd
    \item Use the Stiefel manifold formalism to compute mean/PCA curves
    \item Use the distance (MMD discrepancy test) to mean curve as the decision proxy
\end{itemize}
Strengths: I do think that the area of shape representation using mathematical tools is valuable and therefore I think what the paper proposes would benefit to the community.

%%%%%%%%%%%%%%%%
\begin{point}
I feel that the paper tends to overcomplicate concepts and mix them in a unpractical way. I had a very hard time reading it while the concepts presented stay relatively simple for the aware reader. Some parts feel completely useless or disordered. For example, on page 4, the presentation of the curve representation takes a long paragraph but could be described in a simple sentence. The notation paragraph contains barely any notation used (the curve space $\mathcal{C}_d$ is never defined). A whole paragraph on landmark normalization with a lot of discussion is presented, while the choice seems arbitrary. Also, authors never provide visualisations of what they are doing which require the reader to draw himself some of the concepts to understand them. Overall, I feel that the presentation is the main weakness of the paper which has 20 pages that could be easily reduce to 10 pages of text with illustrations when needed.
\end{point}
\begin{reply}
We have made several structural and visual changes to address these concerns. (i)~The bulk of the SST background (previously spanning much of Section~2) has been relocated to Appendix~\ref{secA0}, and the in-text review (Section~\ref{subsec:SST}) has been shortened and renamed ``Review.'' (ii)~Two new illustrations have been added: Figure~\ref{fig:SRQD} (Section~\ref{subsec:func_approx}) visualizes the SRQD as a weighted polar decomposition, and Figure~\ref{fig:placeholder} (Section~\ref{subsec:mfld_learn}) gives a pipeline diagram tracing curves through each stage of the method. (iii)~Definitions (Definitions~\ref{def:CLO}, \ref{def:PRRTI}, \ref{def:eval_func}) and reformatted lemma/theorem statements make the core objects easier to locate. (iv)~The curve space $\mathcal{C}_d$ is now explicitly defined at Eq.~\eqref{eq:reg_Hilbert_space} in Section~\ref{subsec:preshapes}. A further schematic contrasting nonlinear undulation with linear scale, along with additional concision edits, are in preparation for the next pass.
\end{reply}

\begin{point}
The problem that the paper is trying to solve is not easy to understand. I believe it is to group curves by similarity but then why not using clustering based on the proposed approach.
\end{point}
\begin{reply}
The problem statement is stated explicitly at the top of Section~\ref{sec:intro} (``Problem Statement''): given two images, are the curves in one the same as the curves in the other up to rigid motion, and if not, is the difference due to linear (scale) or nonlinear (undulation) deformations? This is quantitatively instantiated in Section~\ref{subsec:MMD}. Regarding clustering versus hypothesis testing: new blue-text paragraphs in Section~\ref{sec:intro} formalize hypothesis testing as a binary mapping $H:(\mathcal{E},\widehat{\mathcal{E}})\mapsto\{0,1\}$---a \emph{decision} with quantifiable statistical guarantees---whereas clustering aggregates without a formal decision rule. Clustering remains a valuable complementary analysis; we are preparing a demonstration (see R4.3) showing that clustering over the $(\boldsymbol{t},\boldsymbol{\ell})$ features \emph{within} a single image can itself be informative.
\end{reply}

\begin{point}
The choice of evaluation, limited to EBSD feels restricted, while the paper is written in a very generic way. This makes it hard to read for EBSD practicioners as well as for Math Imaging people. An evaluation on similar datasets as in \cite{kaltenmark2017} would be great.
\end{point}
\begin{reply}
Our evaluation is not limited to EBSD. Section~\ref{sec:experiments} (``Segmentation Efficacy'') includes a study over SEM images of lithium-ion battery NMC particles, using morphological segmentation to illustrate the method's applicability to generic grayscale imaging. Additional examples are in preparation subject to coauthor review: a SUVI solar-ultraviolet clustering demonstration (showing clustering over $(\boldsymbol{t},\boldsymbol{\ell})$ \emph{within} a single image) and a pairwise MNIST digit-classification experiment. A curve-comparison evaluation along the lines of~\cite{kaltenmark2017} is also under consideration.
\end{reply}

\begin{point}
An algorithm presenting the pipeline would be very helpful.
\end{point}
\begin{reply}
A pipeline diagram has been added as Figure~\ref{fig:placeholder} in Section~\ref{subsec:mfld_learn}, mapping preshape curves through SRQD/polar decomposition, cyclic Procrustes registration, and separate tangent PCA on $Gr(d,n)$ and $S^d_{++}$ to the normal coordinates $(\boldsymbol{t},\boldsymbol{\ell})$ on which MMD is computed. The complementary Figure~\ref{fig:SRQD} (Section~\ref{subsec:func_approx}) provides a zoomed visualization of the SRQD stage. A companion pseudocode algorithm block is in preparation.
\end{reply}

\begin{point}
The author claims a lot in the introduction the ``interpretability'' of the method, but it's not used otherwise. Learned shape features could produce the same results content. Moreover, none of the AI methods the paper target is (1) cited, (2) compared to.
\end{point}
\begin{reply}
The revision addresses both concerns. (1) A new blue-text paragraph at the end of Section~\ref{sec:experiments} walks through how each row of Truth Table~\ref{tab:logical_conjunction} maps to a distinct physical conclusion in the EBSD experiments: $(1\wedge1)$ = no discrepancy (Figure~\ref{fig:compare_grains_111}); $(1\wedge0)$ = scale-only effect, consistent with visible grain-size differences (Figure~\ref{fig:compare_grains_100}); $(0\wedge1)$ = undulation-only effect from boundary smoothing, below visual resolution (Figure~\ref{fig:compare_grains_010}); $(0\wedge0)$ = both, flagging a presumed faulty measurement (Figure~\ref{fig:compare_grains_000}). The dimensionality-reduction visualization in Figure~\ref{fig:low_dim_grain} further makes these distinctions concrete: varying $\boldsymbol{t}$ alone modulates boundary roughness while varying $\boldsymbol{\ell}$ alone modulates anisotropic size. (2) A qualitative comparison with varifold~\cite{kaltenmark2017,pierson2022}, Sobolev-metric~\cite{bauer2017}, and deep-learning~\cite{hartman2023,hartman2021} approaches has been added as blue text following the contributions list in Section~\ref{sec:intro}. Varifold and Sobolev methods produce a single scalar distance without decomposing \emph{why} two shapes differ; deep-learning approaches require labeled training data and produce opaque latent features. Our framework is unsupervised and produces interpretable features by construction via the separable Truth Table~\ref{tab:logical_conjunction}.
\end{reply}

\begin{point}
Some references on the topic are ignored or wrong:
\begin{itemize}
    \item \cite{kaltenmark2017} ---  should probably in place of [38] (page 7 cited as varifolds)
    \item \cite{pierson2022} --- Represents human shapes as landmarks in RHKS - close to some concepts presented here.
    \item \cite{bauer2017} ---  A numerical framework for Sobolev metrics on the space of curves. Similar to [11] but on curves
    \item \cite{hartman2023} --- A learned feature vector for curves
    \item\cite{hartman2021} --- Similar learning work.
\end{itemize}
\end{point}
\begin{reply}
All five references have been added to the manuscript bibliography with \texttt{sn-aps} formatting. Specifically: \cite{kaltenmark2017} is now cited in Section~\ref{subsec:preshapes} alongside the existing varifold references, correcting the previously incomplete citation; \cite{bauer2017} is cited where Sobolev metrics are discussed in Section~\ref{subsec:preshapes}; and \cite{pierson2022}, \cite{hartman2023}, and \cite{hartman2021} appear together with the blue-text qualitative comparison between our framework and varifold/Sobolev/deep-learning methods in the contributions discussion of Section~\ref{sec:intro}. Duplicate \texttt{\textbackslash bibitem} entries and one malformed DOI were corrected during bibliography cleanup; see CHANGES.md for details.
\end{reply}

\begin{thebibliography}{9}
\bibitem{grey2023sst} Grey, Z. J., Doronina, O. A., \& Glaws, A. (2023). Separable shape tensors for aerodynamic design. Journal of Computational Design and Engineering, 10(1), 468-487.
\bibitem{kaltenmark2017} Kaltenmark, I., Charlier, B., \& Charon, N. (2017). A general framework for curve and surface comparison and registration with oriented varifolds.
\bibitem{pierson2022} Pierson, E., Daoudi, M., \& Arguillere, S. (2022, October). 3d shape sequence of human comparison and classification using current and varifolds. 
\bibitem{bauer2017} Bauer, M., Bruveris, M., Harms, P., \& Møller-Andersen, J. (2017).
\bibitem{hartman2023} Hartman, E., \& Pierson, E. (2023). Varigrad: A novel feature vector architecture for geometric deep learning on unregistered data. 
\bibitem{hartman2021} Hartman, E., Sukurdeep, Y., Charon, N., Klassen, E., \& Bauer, M. (2021). Supervised deep learning of elastic SRV distances on the shape space of curves. 
\end{thebibliography}

\end{document}



